\documentclass[12pt,a4paper]{article}

\usepackage[a4paper,margin=2cm]{geometry}
\usepackage[T2A]{fontenc}
\usepackage[utf8]{inputenc}
\usepackage[russian]{babel}
\usepackage{paratype}
\usepackage{amsmath,amssymb,mathtools}
\usepackage{array,booktabs,longtable}
\usepackage{xcolor}
\usepackage{hyperref}
\usepackage{tikz}

\hypersetup{
  colorlinks=true,
  linkcolor=blue!55!black,
  urlcolor=blue!55!black
}
\setlength{\parindent}{0pt}
\setlength{\parskip}{0.55em}
\emergencystretch=2em
\setlength{\tabcolsep}{3pt}
\renewcommand{\arraystretch}{1.2}

\begin{document}
\hypersetup{pageanchor=false}

\begin{titlepage}
\thispagestyle{empty}
\begin{center}
\vspace*{0.22\textheight}
{\LARGE\bfseries Ревью домашней работы\par}
\vspace{2.2cm}
{\large Автор: Лёвин Артём Александрович\par}
\vspace{0.7cm}
{\large 21 сентября 2026 года\par}
\vfill
\begin{tikzpicture}[x=1cm,y=1cm]
  \draw[blue!35!black, line width=0.8pt]
    (-4,0) .. controls (-2.2,0.55) and (2.2,-0.55) .. (4,0);
\end{tikzpicture}
\end{center}
\end{titlepage}
\hypersetup{pageanchor=true}
\pagenumbering{arabic}

\newpage
\section*{Сводная таблица}
\begin{center}
\small
\begin{tabular}{>{\centering\arraybackslash}p{0.07\linewidth} p{0.19\linewidth} p{0.25\linewidth} p{0.17\linewidth} p{0.18\linewidth}}
\toprule
\textbf{№} & \textbf{Тема} & \textbf{Суть ошибки} & \textbf{Ваш ответ} & \textbf{Правильный ответ} \\
\midrule
\hyperref[task:4]{4} & Геометрические утверждения & Свойство равнобедренного треугольника распространено на любую высоту; верным является утверждение о прямоугольном треугольнике с углами $30^\circ$ и $60^\circ$. & 2 & 3 \\
\addlinespace
\hyperref[task:10]{10} & Медиана набора данных & Для медианы использованы пятый и шестой элементы исходного, неупорядоченного ряда. Сначала данные нужно расположить по возрастанию. & $110$ мм рт. ст. & $112{,}5$ мм рт. ст. \\
\bottomrule
\end{tabular}
\end{center}

\newpage
\thispagestyle{empty}
\vspace*{0.35\textheight}
\begin{center}
{\LARGE\bfseries Разбор ошибок}
\end{center}
\newpage

\phantomsection\label{task:4}
\section*{Задание 4}

\textbf{Текст задания.} Укажите верное утверждение.
\begin{enumerate}
  \item Один из смежных углов всегда тупой.
  \item Любая высота равнобедренного треугольника совпадает с его медианой.
  \item Если угол прямоугольного треугольника равен $60^\circ$, то его гипотенуза в два раза больше одного из катетов.
\end{enumerate}

\textbf{Ваше решение.} Промежуточные рассуждения не записаны; выбран вариант 2.

\textbf{Ваш ответ.} 2.

\textcolor{red}{\textbf{Ошибка:}} выбрано неверное геометрическое утверждение.

\textbf{Где именно возникает ошибка.} Последний достоверно корректный элемент записи --- ответ дан в требуемом формате номером варианта. Первый неверный шаг --- выбор второго утверждения как верного.

\textbf{Почему этот шаг неверен.} Во втором утверждении слово «любая» делает фразу неверной. В равнобедренном треугольнике действительно есть важное совпадение: высота, проведённая из вершины между равными сторонами к основанию, одновременно является медианой и биссектрисой. Но из этого не следует, что каждая высота такого треугольника является медианой. Высоты, проведённые из двух вершин основания, в общем случае не делят противоположные стороны пополам. Поэтому свойство одной специальной высоты нельзя переносить на все три высоты равнобедренного треугольника.

\textcolor{blue!60!black}{\textbf{Ключевая идея:}} при проверке утверждения с словами «всегда», «любая», «каждая» нужно убедиться, что свойство выполняется во всех допустимых случаях. Достаточно одного корректного контрпримера, чтобы такое утверждение оказалось ложным.

\textbf{Исправление.} Проверим варианты по очереди.

Первое утверждение неверно. Смежные углы в сумме дают $180^\circ$, но они могут быть равны по $90^\circ$. Тогда ни один из них не является тупым.

Второе утверждение неверно. В равнобедренном треугольнике медианой обязательно является высота, проведённая к основанию из вершины между равными сторонами. Для двух остальных высот такого общего свойства нет.

Третье утверждение верно. Если один острый угол прямоугольного треугольника равен $60^\circ$, то второй острый угол равен
\[
90^\circ-60^\circ=30^\circ.
\]
В прямоугольном треугольнике с углом $30^\circ$ катет, лежащий напротив этого угла, равен половине гипотенузы. Значит, гипотенуза в два раза больше этого катета.

\textbf{Полное правильное решение.} Из трёх предложенных утверждений только третье выдерживает проверку во всех случаях. Первое опровергается парой смежных прямых углов, второе неверно из-за слова «любая», а третье следует из свойства прямоугольного треугольника с углами $30^\circ$, $60^\circ$ и $90^\circ$.

\textbf{Проверка.} Для третьего утверждения можно использовать стандартное соотношение сторон треугольника с углами $30^\circ$, $60^\circ$, $90^\circ$: сторона напротив угла $30^\circ$ в два раза меньше гипотенузы. Следовательно, гипотенуза действительно в два раза больше одного из катетов.

\textcolor{teal!60!black}{\textbf{Итог:}} верный вариант --- 3.

\textbf{Задача на закрепление.} Укажите верное утверждение:
\begin{enumerate}
  \item Любая медиана равнобедренного треугольника является высотой.
  \item Если один острый угол прямоугольного треугольника равен $30^\circ$, то катет, лежащий напротив него, в два раза меньше гипотенузы.
  \item Сумма двух вертикальных углов всегда равна $180^\circ$.
\end{enumerate}

\newpage
\phantomsection\label{task:10}
\section*{Задание 10}

\textbf{Текст задания.} В таблице даны результаты измерений систолического артериального давления десяти студентов во время дистанционного занятия:
\[
110,\ 125,\ 105,\ 115,\ 120,\ 100,\ 110,\ 105,\ 130,\ 135.
\]
Найдите медиану величины «систолическое артериальное давление» в группе обследованных студентов. Ответ дайте в мм рт. ст.

\textbf{Ваше решение.}
\[
120+100=220,
\]
\[
220:2=110.
\]

\textbf{Ваш ответ.} $110$ мм рт. ст.

\textcolor{red}{\textbf{Ошибка:}} медиана найдена по середине исходной записи, а не по середине упорядоченного ряда.

\textbf{Где именно возникает ошибка.} Последний правильный шаг --- замечено, что при десяти значениях медиану нужно получать как среднее двух центральных элементов. Первый неверный шаг --- в качестве этих элементов выбраны $120$ и $100$, то есть пятое и шестое числа в исходном порядке. Для медианы позиции отсчитываются только после упорядочивания данных.

\textbf{Почему этот шаг неверен.} Медиана определяется положением значений в вариационном ряду, то есть в наборе, расположенном по возрастанию или по убыванию. Исходный порядок наблюдений не несёт информации о том, какие числа являются центральными по величине. Поэтому нельзя брать два числа, просто стоящие посередине напечатанного списка. Сначала все десять результатов нужно упорядочить. Поскольку количество результатов чётное, после упорядочивания центральными будут пятый и шестой элементы.

\textcolor{blue!60!black}{\textbf{Ключевая идея:}} медиана зависит не от места числа в исходной таблице, а от его места после сортировки всего набора.

\textbf{Исправление от последнего правильного шага.} Сохраняем верную идею: значений десять, поэтому после сортировки нужно усреднить два центральных. Расположим данные по возрастанию:
\[
100,\ 105,\ 105,\ 110,\ 110,\ 115,\ 120,\ 125,\ 130,\ 135.
\]
Пятый элемент равен $110$, шестой --- $115$. Поэтому
\[
\frac{110+115}{2}=\frac{225}{2}=112{,}5.
\]

\textbf{Полное правильное решение.} В наборе десять результатов. Упорядочиваем их по возрастанию:
\[
100,\ 105,\ 105,\ 110,\ 110,\ 115,\ 120,\ 125,\ 130,\ 135.
\]
При чётном количестве элементов медиана равна среднему арифметическому двух центральных значений. Здесь это пятое и шестое значения:
\[
110 \quad \text{и} \quad 115.
\]
Следовательно,
\[
\frac{110+115}{2}=112{,}5.
\]

\textbf{Проверка.} После сортировки слева от промежутка между $110$ и $115$ находятся пять значений с учётом центрального левого элемента, а справа --- пять значений с учётом центрального правого элемента. Число $112{,}5$ лежит ровно между двумя центральными значениями, поэтому оно соответствует определению медианы для чётного количества наблюдений.

\textcolor{teal!60!black}{\textbf{Итог:}} медиана равна $112{,}5$ мм рт. ст.

\textbf{Задача на закрепление.} Для набора результатов
\[
132,\ 118,\ 125,\ 110,\ 121,\ 115,\ 129,\ 123
\]
найдите медиану.

\vfill
\begin{center}
\small Лёвин Артём Александрович эксклюзивно для Екатерины Скелли
\end{center}

\end{document}
